\chapter{Conclusiones}
DME responde la pregunta de investigación con una referencia por archivo y destino, condicionada a la ubicación efectiva y al relevo obligatorio de CondorIO. La razón expresa proximidad volumétrica a esa cota; el desglose explica las copias adicionales. Con requerido positivo común, mayor DME equivale a menor observado.

La campaña prospectiva completa las 63 ejecuciones principales previstas en doce condiciones, separadas de doce pilotos válidos y un fallo técnico preservado. Seis repeticiones por condición sintética comprueban tres comportamientos: pipeline de 60 a 50 y 40 MiB con 20 MiB requeridos; archivo compartido de 100 a 80 y 70 MiB con 70 MiB requeridos; particiones distintas en 40 MiB con DME igual a uno. Los cambios se explican por reenvíos y entradas repetidas.

Agrupar las tres etapas conserva retornos intermedios. Omitirlos mediante una intervención separada reduce el observado de 40 a 20 MiB y lleva DME de $1/2$ a $1$, manteniendo el final. En Montage, tres repeticiones por variante en W1 requieren 2\,315\,771 bytes. Los observados al separar, agrupar y omitir además el retorno intermedio son, respectivamente, 3\,375\,611, 2\,845\,691 y 2\,315\,771 bytes. Los finales conservan dimensiones, máscaras y píxeles. Se demuestra exceso evitable en esos dos grafos; el contraejemplo analítico de jobs aislados sin caché impide generalizar esa conclusión.

GL3 y GD4 entregan 40 MiB con requeridos distintos. DME distingue su relación con las obligaciones sin establecer superioridad global. El contraste histórico libre muestra además que una razón mayor puede acompañar más bytes. Para elegir por volumen se compara el observado; para interpretarlo se conservan la referencia y su desglose.

Los tres primeros objetivos quedan respaldados por la caracterización, definición y cálculo implementado. OE4 comprueba correctitud, consistencia, sensibilidad y límites, pero permanece parcial en el costo de captura nativa. La evidencia comprende flujos pequeños, dos operaciones de Montage y un único entorno CondorIO. Las duraciones describen configuraciones completas y no aíslan el efecto temporal de los bytes.
